\begingroup
\renewcommand{\arraystretch}{1.45}
\fancyhead[L]{\small DETERMINATION OF ATTRACTION}
\Needspace{42\baselineskip}
\AttractionArticle{10}

Our formulas can become indeterminate in certain cases, which must be considered separately. First let us examine the case in which the negative roots $-G'$, $-G''$ of the cubic equation become equal, so that, by formulas 18, 19, the coefficients $\gamma'$, $\gamma''$ seem to acquire infinite values, which, however, are in fact indeterminate.

Setting $g=G$ in formula 14, it is evident that, for the two values of $x$, that is, $-G'$ and $-G''$, to become equal, it is necessary that
\[
\begin{aligned}
&AB=0,\quad aa-bb-\frac{aaAA}{aa+G}+\frac{bbBB}{bb+G}=0
\end{aligned}
\]
\par Hence, since $aa-bb$ is by its nature either a positive quantity or $=0$, it readily follows that the following must hold:
\[
\begin{aligned}
&B=0
\end{aligned}
\]
\[
\begin{aligned}
&aa-bb=\frac{aaAA}{aa+G},\quad\text{or}\quad aa+G=\frac{aaAA}{aa-bb}
\end{aligned}
\]
Substituting these values in equation 14 gives
\[
\begin{aligned}
&G'=G''=bb
\end{aligned}
\]
Substituting further the value $x=-bb$ in the cubic equation 13 gives
\[
\begin{aligned}
&(aa-bb)(CC+bb)=bbAA
\end{aligned}
\]
Whenever this conditional equation holds together with the equation $B=0$, the case treated here occurs. \Needspace{31\baselineskip}
And since
\[
\begin{aligned}
&G=\frac{aaAA}{aa-bb}-aa=\frac{aaCC}{bb}
\end{aligned}
\]
\noindent\SourceLeaf{356}
formula 17 gives
\[
\begin{aligned}
&\gamma=\surd{\frac{aabbAA}{(aa-bb)(aaCC+b^4)}}
=\surd{\frac{aaCC+aabb}{aaCC+b^4}}
\end{aligned}
\]
and then formulas 3, 4 give
\[
\begin{aligned}
&\alpha=\frac{\gamma(aa-bb)}{aA}=\frac{\gamma bbA}{a(CC+bb)}
=\surd{\frac{bb(aa-bb)}{aaCC+b^4}}
\\[6pt]
&\qquad{}=\surd{\frac{b^4AA}{(CC+bb)(aaCC+b^4)}}
\end{aligned}
\]
\[
\begin{aligned}
&\beta=0
\end{aligned}
\]
\Needspace{29\baselineskip}
\par The values of the coefficients $\gamma'$, $\gamma''$ given by formulas 18, 19 remain indeterminate in this case, and so do the values of the remaining coefficients $\alpha'$, $\beta'$, $\alpha''$, $\beta''$. Nevertheless, all five of the others can be expressed in terms of one of these coefficients; for example, formula 6 gives
\[
\begin{aligned}
&\alpha'=\frac{\gamma'aA}{aa-bb}
\end{aligned}
\]
and then
\[
\begin{aligned}
&\beta'=\surd{(1-\alpha'\alpha'+\gamma'\gamma')},\\[6pt]
&
\gamma''=\surd{(\gamma\gamma-1-\gamma'\gamma')},\\[6pt]
&
\alpha''=\frac{\gamma''aA}{aa-bb},\\[6pt]
&
\beta''{}=\surd{(1-\alpha''\alpha''+\gamma''\gamma'')}
\end{aligned}
\]
\Needspace{19\baselineskip}
But this is done more elegantly as follows. From
\[
\begin{aligned}
&\gamma\gamma=1+\alpha\alpha,\quad
\alpha\alpha'=\gamma\gamma',\quad
1=\alpha'\alpha'+\beta'\beta'-\gamma'\gamma'
\end{aligned}
\]
it follows that
\[
\begin{aligned}
&\beta'\beta'+\frac{\gamma'\gamma'}{\alpha\alpha}
=1-\alpha'\alpha'+\frac{\gamma\gamma\gamma'\gamma'}{\alpha\alpha}=1
\end{aligned}
\]
\Needspace{32\baselineskip}
\par We may therefore set
\[
\begin{aligned}
&\beta'=\cos f,\quad\gamma'=\alpha\sin f,\quad\alpha'=\gamma\sin f
\end{aligned}
\]
Then, from the formulas
\[
\begin{aligned}
&\varepsilon\alpha''=\beta\gamma'-\gamma\beta',\quad
\varepsilon\beta''=\gamma\alpha'-\alpha\gamma',\quad
\varepsilon\gamma''=\beta\alpha'-\alpha\beta',\quad
\varepsilon\varepsilon=1
\end{aligned}
\]
we find
\[
\begin{aligned}
&\alpha''=-\varepsilon\gamma\cos f,\quad
\beta''=\varepsilon\sin f,\quad
\gamma''=-\varepsilon\alpha\cos f
\end{aligned}
\]
The value of the angle $f$ is arbitrary here, and either $\varepsilon=+1$ or $\varepsilon=-1$ may also be chosen at will.

\Needspace{43\baselineskip}
\AttractionArticle{11}

If $G'$, $G''$ are unequal, the values of the coefficients $\gamma$, $\gamma'$, $\gamma''$ given by formulas 17, 18, 19 cannot be indeterminate. But whenever one of the quantities
\noindent\SourceLeaf{357}
$aa-G'$, $bb-G'$, $aa-G''$, $bb-G''$ vanishes, the value of the coefficient $\alpha'$, $\beta'$, $\alpha''$, $\gamma''$ given by formula 6, 7, 9, 10 respectively seems at first sight to remain indeterminate; a little attention will show, however, that this is otherwise.

Let us suppose, for example, that $aa-G'=0$; then, by equation 18, $\gamma'=0$, and also, by equation 7, $\beta'=0$ (provided that $aa=bb$ does not hold at the same time), whence it is necessary that $\alpha'=\pm1$. If, however, $aa=bb$ at the same time, the formula preceding the sixth in Art. 5 gives $\alpha'A+\beta'B=0$; this equation, combined with $\alpha'\alpha'+\beta'\beta'=1$, yields
\[
\begin{aligned}
&\alpha'=\frac{B}{\surd{(AA+BB)}},\quad
\beta'=\frac{-A}{\surd{(AA+BB)}}
\end{aligned}
\]
\par These expressions evidently cannot be indeterminate unless $A=0$, $B=0$ at the same time; but then we shall return to the case already considered in the preceding article.

\clearpage
\endgroup
